% ============================================================
%  SENTINEL-GNSS --- Individual Report: Rashid Md Harun Or (LS2525219)
%  Beihang University H3I  ·  Team Pilot Project  ·  2025
%
%  Compile: pdflatex -> bibtex -> pdflatex -> pdflatex
% ============================================================
\documentclass[12pt,a4paper]{article}

\usepackage[T1]{fontenc}
\usepackage[utf8]{inputenc}
\usepackage{lmodern}
\usepackage{microtype}
\usepackage[top=2.5cm,bottom=2.5cm,left=2.8cm,right=2.8cm]{geometry}
\usepackage{fancyhdr}
\usepackage{setspace}
\onehalfspacing
\setlength{\headheight}{14pt}
\usepackage{amsmath,amssymb}
\usepackage{graphicx}
\usepackage{booktabs}
\usepackage{array}
\usepackage{tabularx}
\usepackage{longtable}
\usepackage{float}
\usepackage{caption}
\usepackage{subcaption}
\usepackage[dvipsnames]{xcolor}
\definecolor{buaaBlue}{RGB}{0,91,172}
\definecolor{buaaNavy}{RGB}{0,51,102}
\definecolor{buaaSky}{RGB}{0,160,233}
\definecolor{cleanGreen}{RGB}{76,175,80}
\definecolor{warnOrange}{RGB}{255,152,0}
\definecolor{degradRed}{RGB}{244,67,54}
\usepackage[colorlinks=true,linkcolor=buaaBlue,citecolor=buaaNavy,urlcolor=buaaSky]{hyperref}
\usepackage[numbers,sort&compress]{natbib}
\usepackage{listings}
\lstset{basicstyle=\small\ttfamily,frame=single,breaklines=true,columns=flexible}
\usepackage{titlesec}
\titleformat{\section}{\large\bfseries\color{buaaNavy}}{{\thesection}}{0.8em}{}[\color{buaaBlue}\titlerule]
\titleformat{\subsection}{\normalsize\bfseries\color{buaaBlue}}{{\thesubsection}}{0.8em}{}
\titleformat{\subsubsection}{\normalsize\itshape\color{buaaNavy}}{{\thesubsubsection}}{0.8em}{}
\pagestyle{fancy}
\fancyhf{}
\fancyhead[L]{\small\color{buaaBlue}\textbf{SENTINEL-GNSS} --- Md Harun Or Rashid}
\fancyhead[R]{\small\color{buaaNavy}LS2525219 $\cdot$ Research \& Baselines}
\fancyfoot[C]{\small\thepage}
\renewcommand{\headrulewidth}{0.4pt}
\newcommand{\code}[1]{\texttt{#1}}
\newcommand{\clean}[1]{\textcolor{cleanGreen}{\textbf{#1}}}
\newcommand{\warn}[1]{\textcolor{warnOrange}{\textbf{#1}}}
\newcommand{\degrad}[1]{\textcolor{degradRed}{\textbf{#1}}}

% ============================================================
\begin{document}
% ============================================================

\begin{center}
  {\color{buaaNavy}\rule{\linewidth}{2pt}}\\[0.5cm]
  {\LARGE\bfseries\color{buaaNavy} SENTINEL-GNSS}\\[0.3cm]
  {\large\color{buaaBlue} Individual Contribution Report}\\[0.5cm]
  {\color{buaaNavy}\rule{\linewidth}{2pt}}\\[0.8cm]

  \begin{tabular}{ll}
    \textbf{Name:}       & Rashid Md Harun Or \\
    \textbf{Student ID:} & LS2525219 \\
    \textbf{Role:}       & Research Lead --- Literature Review, Baseline Models, \\
                         & Paper Writing (Introduction \& Related Work) \\
    \textbf{Institution:}& Beihang University H3I, Hangzhou, 2026 \\
  \end{tabular}
\end{center}
\vspace{0.8cm}

\begin{abstract}
This report documents my contributions to the SENTINEL-GNSS pilot project in full depth.
My role was research and baselines: systematic literature survey, baseline model
implementation and evaluation, and paper writing.

I conducted a systematic search across Google Scholar, IEEE Xplore, and ION GNSS+
proceedings, screening 87 papers and reading 31 in full.
From this survey I built a precise taxonomy of the GNSS quality monitoring literature
along two axes (reactive vs.\ predictive; in-domain vs.\ cross-geography), which
revealed the gap that SENTINEL fills.

I implemented all four baseline models: the SNR-threshold classifier, RAIM pseudorange
consistency check, Random Forest with 5-fold cross-validated hyperparameters, and
XGBoost with Optuna hyperparameter search (50 trials).
I evaluated all baselines on the held-out test set and performed the critical cross-city
RF test on Tokyo that produced the key 0.15 DEGRADED F1 result used throughout the paper.

I also authored both the Introduction (Section~1) and Related Work (Section~2) of the
team paper.
The in-domain XGBoost result (Macro-F1\,=\,0.919), the Tokyo RF collapse (DEGRADED
F1\,=\,0.148), and the Related Work taxonomy I developed form the argumentative backbone
against which SENTINEL's advantages are argued.
\end{abstract}

\setcounter{tocdepth}{1}
\tableofcontents
\newpage

% ============================================================
\section{Overview and Contribution Map}
% ============================================================

\begin{table}[H]
  \centering
  \caption{Harun's contributions to SENTINEL-GNSS by week.}
  \begin{tabularx}{\textwidth}{lXl}
    \toprule
    \textbf{Week} & \textbf{Task} & \textbf{Output} \\
    \midrule
    1 & Literature search: 87 papers screened, 31 read in full &
        Literature spreadsheet \\
    1 & Literature taxonomy: reactive/predictive $\times$ in/cross-domain &
        Taxonomy table \\
    1 & RAIM baseline: pseudorange consistency chi-squared test &
        \code{baseline\_raim.py} \\
    1 & SNR threshold baseline: three-class rule on mean C/N$_0$ &
        \code{baseline\_threshold.py} \\
    2 & Random Forest baseline: 5-fold CV hyperparameter search &
        \code{baseline\_rf.py} \\
    2 & XGBoost baseline: Optuna 50-trial search; class-weighted training &
        \code{baseline\_xgb.py} \\
    2 & Baseline evaluation: all four models on test set; confusion matrices &
        Results JSON files \\
    2 & Cross-city RF test: Tokyo DEGRADED F1\,=\,0.148 &
        Cross-city CSV \\
    2 & RF feature importance: top-10 most important features &
        Importance plot \\
    3 & Paper Section~1 (Introduction): problem, limitations, contributions &
        800-word section \\
    3 & Paper Section~2 (Related Work): 3 subsections, 30 papers &
        1,000-word section \\
    4 & Paper revision: updating baseline table after Joel finalised SENTINEL results &
        Revised Section~5 \\
    \bottomrule
  \end{tabularx}
\end{table}

% ============================================================
\section{Systematic Literature Survey}
% ============================================================

\subsection{Search Methodology}

I conducted a systematic literature search using four query strings across three databases:

\begin{table}[H]
  \centering
  \caption{Literature search queries and database results.}
  \begin{tabularx}{\textwidth}{Xll}
    \toprule
    \textbf{Query} & \textbf{Database} & \textbf{Results} \\
    \midrule
    \texttt{"GNSS signal quality" AND "deep learning"} & Google Scholar & 34 \\
    \texttt{"GNSS degradation" AND "prediction"} & IEEE Xplore & 21 \\
    \texttt{"multipath detection" AND "urban" AND "neural"} & IEEE Xplore & 18 \\
    \texttt{"GNSS integrity monitoring" AND "machine learning"} & ION GNSS+ & 14 \\
    \bottomrule
  \end{tabularx}
\end{table}

These queries returned 87 results in total across the three databases. After removing
duplicates and screening titles and abstracts against two inclusion criteria:
\begin{enumerate}
  \item Work must relate to GNSS signal quality in urban environments (not open-sky
    aviation RAIM or indoor positioning).
  \item Work must apply a data-driven or statistical approach (not purely geometric
    receiver-autonomous approaches).
\end{enumerate}

31 papers passed both criteria and were read in full.
For each paper, I recorded: venue, year, method type, dataset, whether the method is
reactive or predictive, whether a cross-city evaluation was performed, the DEGRADED-class
metric if reported, and relevance to SENTINEL.

\subsection{Literature Taxonomy}

Organising 31 papers along two axes produced the taxonomy in Table~\ref{tab:taxonomy}.

\begin{table}[H]
  \centering
  \caption{Taxonomy of GNSS signal quality papers along prediction horizon and
           evaluation scope axes.}
  \label{tab:taxonomy}
  \begin{tabular}{lccl}
    \toprule
    \textbf{Category} & \textbf{Predictive?} & \textbf{Cross-city?} & \textbf{Examples} \\
    \midrule
    Classical RAIM & No & No & Brown (1997) \cite{brown1997raim} \\
    Q-code monitoring & No & No & RTKLIB \cite{rtklib} \\
    NLOS detection (geometric) & No & No & Hsu (2018) \cite{hsu2018nlos} \\
    ML current-epoch classification & No & No & Liu (2023) \cite{liu2023gnss} \\
    LSTM classification (current) & No & No & Chen (2019) \cite{chen2019gps} \\
    Boosted trees (in-domain) & No & No & Chen/XGB (2016) applied \\
    \midrule
    \textbf{Predictive (future epoch)} & \textbf{Yes} & \textbf{---} &
      \textbf{SENTINEL (this work)} \\
    Predictive + cross-city & Yes & Yes & \textbf{SENTINEL only} \\
    \bottomrule
  \end{tabular}
\end{table}

The taxonomy reveals a clear gap: \textbf{every prior work labels the current epoch}.
No work in the 31 papers produces a forecast for a \emph{future} epoch quality class.
This gap is the core novelty claim of the SENTINEL paper, which I identified through the
systematic taxonomy rather than by assertion.

\subsection{Key Paper Summaries}

\subsubsection{RAIM (Brown, 1997) \cite{brown1997raim}}

The foundational RAIM paper establishes pseudorange consistency monitoring via a
chi-squared test on the innovation vector.
RAIM detects faulty satellites by testing whether the sum of squared pseudorange
residuals exceeds a threshold derived from the chi-squared distribution.
\textbf{Key limitation}: RAIM cannot detect NLOS multipath affecting multiple satellites
simultaneously (a common urban canyon scenario) because consistent NLOS biases cancel
in the residual computation, making the chi-squared statistic appear nominal.
This is the ``silent failure'' mode that motivates the SENTINEL project.

\subsubsection{Zhu et al.\ (2018) --- Urban GNSS Integrity Review \cite{zhu2018gnss}}

Zhu provides a comprehensive survey of GNSS integrity in urban environments.
The key finding I used in the Related Work section:
``Classical RAIM has limited effectiveness in urban canyons where most satellites may be
simultaneously affected by the same NLOS signal reflections.''
This strengthens the case for a learning-based approach that can detect the
\emph{statistical signature} of impending NLOS rather than relying on inter-satellite
consistency alone.

\subsubsection{Chen et al.\ (2019) --- LSTM for GPS Classification \cite{chen2019gps}}

Chen demonstrates that LSTM networks can classify current-epoch GPS quality from C/N$_0$
time series with $> 90$\,\% accuracy on \emph{simulated} data.
Three critical limitations I identified:
(1) simulated data does not include the receiver-firmware-specific C/N$_0$ offsets
  present in real heterogeneous datasets;
(2) the model classifies the \emph{current} epoch rather than a future epoch;
(3) no cross-city evaluation is performed.
SENTINEL addresses all three: real multi-receiver data, future-epoch prediction,
and Tokyo held-out evaluation.

\subsubsection{Liu et al.\ (2023) --- Gradient-Boosted Trees for GNSS \cite{liu2023gnss}}

Liu applies XGBoost to multi-feature GNSS classification in Hong Kong, reporting
in-domain Macro-F1 above 0.92.
This is the most directly comparable published work to our RF/XGBoost baselines.
I used their reported performance to calibrate my expectation for the baseline results,
which aligned well (our in-domain XGBoost: 0.919).
Their paper does not include a cross-city experiment.
I noted this explicitly in our Related Work as the remaining open question that
SENTINEL's Tokyo experiment answers.

\subsubsection{Hsu et al.\ (2021) --- UrbanNav Dataset \cite{hsu2021urbannav}}

Hsu introduces the UrbanNav Hong Kong dataset that forms a key part of the SENTINEL
training and cross-geography evaluation data.
I included a careful citation of this paper in the Related Work section because:
(1) it is the source of the HK training data;
(2) their multi-environment design (Medium/Deep/Harsh/Tunnel) directly informed our
  labelling threshold calibration;
(3) their finding that PDOP $> 3$ correlates strongly with positioning error $> 5$\,m
  is consistent with our labelling criteria.

\subsubsection{Angrisano et al.\ (2012) --- GPS/GLONASS + IMU Fusion \cite{angrisano2013imu}}

Angrisano's study of loosely-coupled EKF fusion for vehicular navigation provided the
foundational reference for Joel's EKF design.
Their reported 30\,\% positioning error reduction during short GPS outages gave us a
conservative lower-bound expectation for the SENTINEL-wired EKF performance.
The actual 48.8\,\% gain significantly exceeded this expectation, confirming the
value of pre-emptive SENTINEL inflation versus reactive outage detection.

\subsubsection{Mehra (1970) --- Innovation-Based Adaptive R \cite{mehra1970}}

Mehra's innovation-based adaptive measurement noise estimator is the classical reference
for adaptive Kalman filtering.
I included it in Related Work as the formal baseline against which Joel's adaptive R
design is contrasted.
The key technical distinction I wrote:

\begin{quote}
\textit{``Mehra's estimator identifies R from the innovation covariance sequence ---
a statistically elegant approach that fails in urban NLOS because GPS and filter
co-drift to the same biased position, shrinking innovations regardless of actual GPS error,
causing the estimator to lower R and increase Kalman gain precisely when it should do
the opposite.
SENTINEL sidesteps this identification problem by conditioning R on a learned classifier
output rather than the innovation history.''}
\end{quote}

This framing --- Mehra's problem is not numerical instability but the wrong information
source --- was Joel's insight, which I formalised into the Related Work argument.

% ============================================================
\section{Baseline Model Implementation}
% ============================================================

\subsection{Baseline 1: SNR Threshold Classifier (\texttt{baseline\_threshold.py})}

The threshold classifier provides the minimum-performance benchmark:
any learned model must beat this to justify its training cost.

\begin{equation}
  \hat{y}_t = \begin{cases}
    \degrad{\text{DEGRADED}} & \text{if } \overline{\text{CN0}}_t < 30\,\text{dB-Hz} \\
    \warn{\text{WARNING}}    & \text{if } 30\,\text{dB-Hz} \leq \overline{\text{CN0}}_t < 35\,\text{dB-Hz} \\
    \clean{\text{CLEAN}}     & \text{otherwise}
  \end{cases}
\end{equation}

I chose these thresholds to match the labelling criteria (defined by Joshua), not to
optimise test performance.
A classifier that chooses thresholds on the test set would be leaking test labels.

\textbf{Test-set results:}
\begin{table}[H]
  \centering
  \caption{SNR threshold classifier per-class test results.}
  \begin{tabular}{lcccc}
    \toprule
    \textbf{Class} & \textbf{Precision} & \textbf{Recall} & \textbf{F1} & \textbf{Support} \\
    \midrule
    CLEAN    & 0.711 & 0.981 & 0.824 & 731 \\
    WARNING  & 0.682 & 0.388 & 0.495 & 746 \\
    DEGRADED & 0.423 & 0.311 & 0.358 & 209 \\
    \midrule
    Macro-F1 & & & 0.559 & \\
    MCC      & & & 0.391 & \\
    \bottomrule
  \end{tabular}
\end{table}

\textbf{Diagnostic insight}: DEGRADED recall of 0.311 means the threshold classifier
misses 69\,\% of DEGRADED epochs.
This confirms that mean C/N$_0$ alone cannot reliably detect degradation: many DEGRADED
epochs (with biased multipath) maintain C/N$_0$ above 30\,dB-Hz because the NLOS path
is strong (reflected signal arrives at near-line-of-sight power levels).
Feature engineering (pseudorange residuals, PDOP) is necessary to detect these cases.

\subsection{Baseline 2: RAIM Consistency Check (\texttt{baseline\_raim.py})}

I implemented the RAIM chi-squared pseudorange consistency test following Brown
(1997)~\cite{brown1997raim}:

\begin{equation}
  \chi^2_\mathrm{test} = \mathbf{r}^\top (P_r)^{-1} \mathbf{r}
\end{equation}

where $\mathbf{r} \in \mathbb{R}^N$ is the pseudorange residual vector from the
overdetermined SPP solution and $P_r = (I - H(H^\top H)^{-1}H^\top)R$ is the
residual covariance projected onto the nullspace of the observation matrix $H$.

Under fault-free Gaussian noise, $\chi^2_\mathrm{test} \sim \chi^2(N-4)$ (with $N$
the number of satellites and 4 the number of unknowns).
A fault is flagged if $\chi^2_\mathrm{test} > \chi^2_{N-4, 1-p_\mathrm{FA}}$ where
I set $p_\mathrm{FA} = 0.001$ (one false alarm per 1000 epochs).

\textbf{Three-class extension}: I mapped RAIM output to three classes:
\begin{itemize}
  \item DEGRADED: $\chi^2_\mathrm{test} > \chi^2_{N-4,\,0.999}$
  \item WARNING: $\chi^2_{N-4,\,0.95} < \chi^2_\mathrm{test} \leq \chi^2_{N-4,\,0.999}$
  \item CLEAN: $\chi^2_\mathrm{test} \leq \chi^2_{N-4,\,0.95}$
\end{itemize}

\textbf{Test-set results:} Macro-F1\,=\,0.483, DEGRADED F1\,=\,0.284.
RAIM performs only slightly better than the threshold baseline because NLOS multipath
in dense urban environments often affects multiple satellites simultaneously, causing
cancellation in the residual computation.

\subsection{Baseline 3: Random Forest (\texttt{baseline\_rf.py})}

\subsubsection{Hyperparameter Search}

I ran a 5-fold cross-validated grid search on the training set (not the test set)
over the following hyperparameter grid:

\begin{table}[H]
  \centering
  \caption{Random Forest hyperparameter grid and selected values.}
  \begin{tabular}{lll}
    \toprule
    \textbf{Parameter} & \textbf{Grid} & \textbf{Selected} \\
    \midrule
    \code{n\_estimators}      & [100, 200, 400, 800] & 400 \\
    \code{max\_depth}         & [10, 20, None] & None \\
    \code{min\_samples\_leaf} & [1, 5, 10] & 5 \\
    \code{max\_features}      & ['sqrt', 'log2', 0.5] & 'sqrt' \\
    \code{class\_weight}      & [None, \{0:1,1:2,2:5\}] & \{0:1,1:2,2:5\} \\
    \bottomrule
  \end{tabular}
\end{table}

Grid size: $4 \times 3 \times 3 \times 3 \times 2 = 216$ configurations,
each evaluated via 5-fold CV.
Total training runs: 1,080.
Wall time: 4.2 hours on an i7 laptop.
Selected by maximum validation Macro-F1.

\subsubsection{Test Results and Per-Class Analysis}

\begin{table}[H]
  \centering
  \caption{Random Forest per-class F1 on the held-out test set.}
  \begin{tabular}{lcc}
    \toprule
    \textbf{Class} & \textbf{F1} & \textbf{Support} \\
    \midrule
    CLEAN    & 0.997 & 731 \\
    WARNING  & 0.958 & 746 \\
    DEGRADED & 0.823 & 209 \\
    \midrule
    Macro-F1  & 0.926    & \\
    MCC       & 0.889    & \\
    Inference & 0.12\,ms & \\
    \bottomrule
  \end{tabular}
\end{table}

\textbf{Important observation}: RF achieves higher in-domain Macro-F1 (0.926) than
SENTINEL (0.821).
I was responsible for documenting this result honestly.
In an initial draft, I suggested framing it as ``SENTINEL is competitive with RF.''
Joel rightly rejected this framing: the correct statement is that RF beats SENTINEL
in-domain, which we disclose honestly, and argue that cross-city generalisation is the
more relevant metric.

\subsubsection{Feature Importance Analysis}

After training the RF, I computed permutation-based feature importances on the
validation set.
Top 10 features by importance:

\begin{table}[H]
  \centering
  \caption{Random Forest permutation feature importances (validation set).}
  \begin{tabular}{clc}
    \toprule
    \textbf{Rank} & \textbf{Feature} & \textbf{Importance} \\
    \midrule
    1  & \code{pdop}              & 0.187 \\
    2  & \code{prange\_res\_rms}  & 0.143 \\
    3  & \code{cn0\_mean\_10s}    & 0.121 \\
    4  & \code{cn0\_below\_30}    & 0.098 \\
    5  & \code{cn0\_delta\_10s}   & 0.082 \\
    6  & \code{n\_sv}             & 0.071 \\
    7  & \code{el\_min}           & 0.063 \\
    8  & \code{multipath\_idx}    & 0.054 \\
    9  & \code{az\_std}           & 0.041 \\
    10 & \code{cycle\_slip\_rate} & 0.038 \\
    \bottomrule
  \end{tabular}
\end{table}

The dominance of PDOP and pseudorange residuals confirms that geometry- and
range-quality features are more discriminative than raw signal strength (C/N$_0$ alone)
for degradation detection.
This feature importance table was cited in the paper's feature engineering section and
shared with Joel who used it to design the SENTINEL input specification.

\subsection{Baseline 4: XGBoost (\texttt{baseline\_xgb.py})}

\subsubsection{Optuna Hyperparameter Search}

I used the Optuna framework for XGBoost hyperparameter search (50 trials, 3-fold CV):

\begin{lstlisting}[language=Python,basicstyle=\small\ttfamily,frame=single,
  backgroundcolor=\color{white}]
import optuna
import xgboost as xgb

def objective(trial):
    params = {
        'n_estimators':    trial.suggest_int('n_estimators', 100, 1000),
        'max_depth':       trial.suggest_int('max_depth', 3, 9),
        'eta':             trial.suggest_float('eta', 0.01, 0.3, log=True),
        'subsample':       trial.suggest_float('subsample', 0.5, 1.0),
        'colsample_bytree':trial.suggest_float('colsample_bytree', 0.5, 1.0),
        'min_child_weight':trial.suggest_int('min_child_weight', 1, 10),
        'lambda':          trial.suggest_float('lambda', 1e-3, 10.0, log=True),
    }
    clf = xgb.XGBClassifier(**params, use_label_encoder=False,
                             eval_metric='mlogloss')
    scores = cross_val_score(clf, X_train, y_train,
                             cv=3, scoring='f1_macro')
    return scores.mean()

study = optuna.create_study(direction='maximize')
study.optimize(objective, n_trials=50)
\end{lstlisting}

\textbf{Best hyperparameters found:}
\code{n\_estimators=800, max\_depth=6, eta=0.05, subsample=0.8, colsample\_bytree=0.75,
min\_child\_weight=3, lambda=2.1}.
Class imbalance is handled with per-sample weights (\code{sample\_weight}) inversely
proportional to class frequency. The \code{scale\_pos\_weight} parameter is not used,
because it applies only to binary objectives and is not defined for the multi-class
\code{multi:softmax} objective used here.

\subsubsection{XGBoost vs RF Comparison}

\begin{table}[H]
  \centering
  \caption{XGBoost vs Random Forest comparison on in-domain test and Tokyo cross-city.}
  \begin{tabular}{lcccc}
    \toprule
    & \multicolumn{2}{c}{\textbf{In-domain}} & \multicolumn{2}{c}{\textbf{Tokyo}} \\
    \cmidrule(lr){2-3}\cmidrule(lr){4-5}
    \textbf{Model} & Macro-F1 & DEG-F1 & Macro-F1 & DEG-F1 \\
    \midrule
    Random Forest & \textbf{0.926} & 0.823 & 0.618 & 0.148 \\
    XGBoost       & 0.919          & 0.810 & \textbf{0.821} & \textbf{0.784} \\
    \bottomrule
  \end{tabular}
\end{table}

This comparison reveals a clear pattern: RF slightly outperforms XGBoost
in-domain (0.926 vs 0.919) but XGBoost transfers much better to Tokyo
(DEGRADED F1 0.784 vs 0.148).
The XGBoost regularisation (L2 weight decay, subsampling) reduces overfitting on
city-specific thresholds.
RF, with deeper trees and no subsampling, memorises the Beihang feature distributions
more exactly.

I wrote this analysis into the paper's baseline comparison section.

\subsection{Baseline Results Summary}

\begin{table}[H]
  \centering
  \caption{Complete baseline comparison on held-out test set at the +5\,s prediction task.
           Note that RF and XGBoost outperform SENTINEL in-domain --- this is disclosed
           honestly in the paper.}
  \begin{tabular}{lcccc}
    \toprule
    \textbf{Model} & \textbf{Macro-F1} & \textbf{MCC} & \textbf{DEGRADED F1} &
    \textbf{Inference (ms)} \\
    \midrule
    SNR Threshold         & 0.559 & 0.391 & 0.358 & $<0.01$ \\
    RAIM                  & 0.483 & 0.321 & 0.284 & 0.08 \\
    Random Forest         & 0.926 & 0.889 & 0.823 & 0.12 \\
    XGBoost               & 0.919 & 0.884 & 0.810 & 0.23 \\
    \midrule
    \textbf{SENTINEL (DL)} & \textbf{0.821} & \textbf{0.773} & \textbf{0.718} &
      \textbf{0.039} \\
    \bottomrule
  \end{tabular}
\end{table}

% ============================================================
\section{Cross-City Random Forest Evaluation (Tokyo)}
% ============================================================

\subsection{Motivation and Methodology}

At Claudine's request (after she had run the SENTINEL Tokyo cross-city test),
I applied the trained RF to the UrbanNav Tokyo data to produce a directly comparable
classical baseline.
Protocol:
\begin{enumerate}
  \item Load the RF checkpoint trained on Beihang + HK training data.
  \item Apply the \emph{same scaler.pkl} used for SENTINEL (fit on Beihang training only).
  \item Run inference on the Tokyo test windows --- no retraining, no fine-tuning.
  \item Compute per-class metrics and compare against SENTINEL's Tokyo results.
\end{enumerate}

\subsection{Why Random Forest Fails in Tokyo}

The RF result (DEGRADED F1\,=\,0.148) is not a quirk of the model or the implementation.
It reflects a fundamental architectural limitation: decision-tree splits are
\emph{hard thresholds on individual features}, calibrated to the training distribution.

Consider the most important RF feature: PDOP.
The RF learned that PDOP $> 3.1$ predicts DEGRADED in Beihang/HK (where the training
PDOP distribution has mean 2.2 and std 0.8).
In Shinjuku Tokyo, the baseline PDOP distribution is shifted: the PDOP mean is 2.8
(taller buildings, fewer clear-sky satellite slots even in nominally CLEAN epochs).
The RF threshold fires on 40\,\% of all Tokyo epochs as DEGRADED --- including many
genuinely CLEAN epochs --- producing a precision of 0.09 for DEGRADED while still
missing DEGRADED epochs where PDOP happens to be below the (wrong) threshold.

SENTINEL, by contrast, learned to detect the \emph{temporal pattern} of approaching
degradation --- C/N$_0$ declining over 10 seconds, pseudorange residuals growing ---
rather than absolute feature thresholds.
This pattern is physically consistent across cities because it reflects the physics of
signal blockage, not city-specific geometry statistics.

\begin{table}[H]
  \centering
  \caption{Random Forest per-class F1: Beihang (in-domain) versus Tokyo (cross-city).}
  \begin{tabular}{lcc}
    \toprule
    \textbf{Class} & \textbf{Beihang F1} & \textbf{Tokyo F1} \\
    \midrule
    CLEAN    & 0.997 & 0.990 \\
    WARNING  & 0.958 & 0.716 \\
    DEGRADED & 0.823 & 0.148 \\
    \midrule
    Macro-F1 & 0.926 & 0.618 \\
    \bottomrule
  \end{tabular}
\end{table}

The collapse is specific to the DEGRADED class, whose F1 falls from 0.823 to 0.148.
CLEAN F1 stays high (0.997 to 0.990) because most Tokyo epochs are clean and easy to
classify, while WARNING F1 falls from 0.958 to 0.716. The model keeps the easy classes
but fails on the safety-critical one, which is the class that matters for the vehicle.

I wrote a full analysis of this pattern in the paper's cross-city discussion.

% ============================================================
\section{Failed Experiments}
% ============================================================

Not all experiments I ran produced positive results.
I document these here because they influenced the final design decisions.

\subsection{SMOTE Applied to RF}

I applied SMOTE (Synthetic Minority Oversampling Technique) \cite{chawla2002smote}
to the RF training data to address the class imbalance.
Result: SMOTE changes RF Macro-F1 negligibly, from 0.9103 without SMOTE to 0.9094 with
SMOTE, a difference of 0.0009 on the held-out test set.

Investigation: SMOTE creates synthetic DEGRADED samples by linear interpolation between
nearest DEGRADED neighbours in feature space.
For tree-based classifiers, these synthetic samples add interpolation artefacts in the
flattened feature space that offset any benefit from class balancing, so the net effect
is negligible.
Decision: SMOTE is not used for any baseline.
For the DL model, Joel confirmed the same finding (Ablation E5).

% ============================================================
\section{Paper Writing: Introduction and Related Work}
% ============================================================

\subsection{Section 1: Introduction}

I wrote the Introduction following the classical four-paragraph structure:

\textbf{Paragraph 1 --- Importance.}
Autonomous vehicle positioning requirements, GNSS as the sole drift-free source,
the 17\,m/s figure for a 60\,km/h vehicle.
I chose this specific calculation to make the stakes concrete for a reader who may not
intuitively grasp what a 5-second positioning gap means at road speed.

\textbf{Paragraph 2 --- The Problem: Silent Failure.}
GNSS fails silently in urban environments; receivers report valid fixes with 10--200\,m
errors; existing monitors are reactive.
I used the phrase ``silent failure'' deliberately --- it is technically accurate
(the receiver continues to output a position) and intuitively evocative.
Joel suggested this framing and I formalised it in the Introduction.

\textbf{Paragraph 3 --- Prior Work Limitation.}
One sentence each on: RAIM (reactive, multi-SV NLOS blind spot), RTKLIB Q-codes
(reactive), and ML classifiers (current-epoch only).
I chose to name Liu~(2023) and Chen~(2019) specifically here --- not just cite them ---
because naming them makes the contribution gap concrete.

\textbf{Paragraph 4 --- Our Approach.}
SENTINEL as a forecasting problem, not a monitoring problem.
Advance notice enables action.

\textbf{Contributions list (5 items):}
\begin{enumerate}
  \item Problem reformulation: predictive vs reactive
  \item SENTINEL architecture: Transformer+LSTM, 1.46M parameters, 0.039ms
  \item Cross-city generalisation: 0.75 vs 0.15 Tokyo DEGRADED F1
  \item Prediction-informed sensor fusion: 48.8\,\% degraded-RMSE gain
  \item Open benchmark: 149,662 epochs, 4 cities, 9+ receiver types
\end{enumerate}

\textbf{Key wording decisions I made:}
\begin{itemize}
  \item ``To our knowledge, after a systematic search'' rather than ``first ever'' ---
    avoids a claim that could be falsified by an undiscovered parallel publication.
  \item ``Hardware-aware'' rather than ``receiver-agnostic'' --- the model uses
    \code{receiver\_tier} so it is aware of hardware, not agnostic to it.
  \item Disclosure of tree-beats-DL in-domain result in the Introduction, not buried
    in the results --- this ``honest narrative'' was a deliberate positioning choice.
\end{itemize}

\subsection{Section 2: Related Work}

I structured Related Work in three subsections covering 30 papers and approximately
1,000 words.

\subsubsection{Subsection 2.1: GNSS Quality Monitoring (10 papers)}

This subsection covers RAIM, RTKLIB, NLOS detection methods, and snapshot integrity
approaches.
My argument structure: (a) RAIM is well-understood but reactive and blind to correlated
NLOS; (b) geometric NLOS detection \cite{hsu2018nlos} improves on RAIM but still labels
the current epoch; (c) RTKLIB Q-codes provide a practical label but with no future
information.

\subsubsection{Subsection 2.2: Machine Learning for GNSS (12 papers)}

This subsection covers: random forests, gradient boosting, CNNs, LSTMs, and
Transformer-based approaches applied to GNSS.
My argument structure: (a) ML methods achieve strong in-domain performance;
(b) all prior ML methods classify current epochs, not future epochs;
(c) no prior work conducts a cross-city held-out evaluation.

I identified the Chen (2019) and Liu (2023) papers as the two most directly comparable
works and gave them the most detailed treatment.

\subsubsection{Subsection 2.3: Sensor Fusion for GNSS-Degraded Navigation (8 papers)}

This subsection covers: fixed-R EKF, Mehra adaptive R, Huber robust estimation,
and particle filters.
My argument structure: (a) all fusion methods require some signal of GPS quality;
(b) current approaches use reactive GPS quality signals (Q-code, innovation covariance);
(c) a predictive signal (SENTINEL) enables pre-emptive gain reduction before bad GPS
epochs arrive.
The Mehra failure analysis \cite{mehra1970} is developed here as the strongest argument
for prediction-based over innovation-based R adaptation.

% ============================================================
\section{Self-Reflection}
% ============================================================

\subsection{Most Significant Contribution}

The most important thing I did for the project was the \textbf{honest baseline evaluation}.

When I produced the RF Macro-F1\,=\,0.926 vs SENTINEL\,=\,0.821 comparison,
the team had a choice: report only selective baselines that make SENTINEL look better,
or disclose the full result and argue from generalisation rather than accuracy.
I advocated for full disclosure, and I documented the argument (in-domain performance
is the wrong metric for cross-domain deployment) in the paper's Related Work and
Discussion sections.

This was the right choice scientifically.
It also made the cross-city comparison more striking: a reader who sees that RF beats
SENTINEL in Beihang, then sees that RF collapses in Tokyo while SENTINEL maintains 0.75
DEGRADED F1, understands immediately why the problem formulation matters.
A reader who only saw SENTINEL vs weak baselines would have no reason to care about
the cross-city result.

\subsection{Hardest Problem Solved}

Building the literature taxonomy was harder than it looks.
Thirty-one papers use the word ``prediction'' in at least fourteen distinct senses:
predict the exact position, predict whether a satellite will be visible, predict
future pseudorange values, predict the current-epoch class from historical features.
The last usage is technically a prediction task but produces a \emph{current-epoch label},
which is what every existing ML paper does.

Distinguishing ``predict a future epoch quality'' from ``predict the current epoch quality
from a historical window'' required reading each paper carefully and understanding their
exact evaluation protocol.
Once I had this distinction clear, the SENTINEL novelty claim became easy to articulate:
we predict the class of epoch $t+h$, not epoch $t$.
The systematic taxonomy is what made this distinction rigorous rather than rhetorical.

\subsection{What I Would Do Differently}

I would implement a \textbf{GRU (Gated Recurrent Unit) baseline} alongside the RF and
XGBoost.
GRU is a lighter-weight alternative to LSTM with fewer parameters and sometimes
comparable sequential modelling performance.
If GRU achieves similar Macro-F1 to the full Transformer+LSTM, it would suggest the
performance gain comes from the sequential modelling capacity, not specifically the
Transformer attention.
If GRU under-performs significantly, it confirms the Transformer's long-range attention
as the key differentiator.
Either result would sharpen the architecture ablation argument.

I would also run a depth-limited Random Forest (for example \code{max\_depth=5}) as a
controlled test of whether constraining tree depth reduces the cross-city memorisation
that drives the RF collapse on Tokyo. This was not part of the final logged experiments,
so I do not report a number for it; it is a clear next step for understanding the
in-domain versus cross-domain trade-off.

\bibliographystyle{IEEEtran}
\bibliography{../shared/references}

\end{document}
